% 原书第 1—5 页；章首导言、第 1.1—1.2 节。
\chapter{统计模型与因果模型}
\label{ch:statistical-causal-models}

借助\textbf{统计学习（statistical learning）}，我们试图从\textbf{观测数据（observational data）}中推断随机变量之间相依关系的性质。例如，根据两个随机变量的一组联合观测样本，我们可以构造一个预测器：当获得其中一个变量的新取值时，用它来较好地估计另一个变量。这类预测所依据的理论已经相当成熟。尽管这些理论通常针对较为简单的情形，却已为我们理解如何从数据中学习提供了深刻认识。

本章将介绍其中一些认识，主要有两个原因。第一，它们有助于我们体会\textbf{因果推断（causal inference）}究竟难在哪里。在因果推断中，基础模型不再是随机变量的某个固定联合分布，而是一种能够导出多个联合分布的结构。第二，尽管关于因果估计的有限样本结果还比较少，我们仍应记住：进入更复杂的因果问题后，基本的统计估计问题并不会消失。与纯粹的统计学习中不会出现的那些因果难题相比，这些统计问题也许显得次要，但仍然需要处理。在这些准备工作的基础上，本章还将借助两个例子，循序渐进地介绍因果性的基本概念；其中一个例子是机器学习领域所熟悉的。

\section{概率论与统计学}
\label{sec:probability-statistics}

概率论与统计学以\textbf{随机试验（random experiment）}或\textbf{概率空间（probability space）} $(\Omega,\mathcal{F},P)$ 为基本模型。其中，$\Omega$ 是包含所有可能结果的集合，$\mathcal{F}$ 是由事件 $A\subseteq\Omega$ 构成的集合族，而 $P$ 是为每个事件赋予概率的测度。给定这一数学结构，概率论使我们能够对随机试验的结果进行推理。统计学习所处理的则基本上是相反的问题：已知试验结果，希望据此推断其背后数学结构的性质。例如，假设我们已经观测到数据
\begin{equation}
 (x_1,y_1),\ldots,(x_n,y_n),
 \tag{1.1}
\end{equation}
其中，$x_i\in\mathcal{X}$ 是\textbf{输入（input）}，有时也称为\textbf{协变量（covariate）}或样例（case）；$y_i\in\mathcal{Y}$ 是\textbf{输出（output）}，有时也称为目标值（target）或\textbf{标签（label）}。我们可以假设，各个 $(x_i,y_i)$（$i=1,\ldots,n$）都由同一个未知随机试验独立生成。更准确地说，这一模型假定，观测值 $(x_1,y_1),\ldots,(x_n,y_n)$ 是随机变量 $(X_1,Y_1),\ldots,(X_n,Y_n)$ 的实现值；这些随机变量\textbf{独立同分布（independent and identically distributed，i.i.d.）}，共同服从\textbf{联合分布（joint distribution）} $P_{X,Y}$。这里，$X$ 和 $Y$ 分别是取值于度量空间 $\mathcal{X}$ 和 $\mathcal{Y}$ 的随机变量。\footnote{随机变量 $X$ 是一个可测函数 $\Omega\to\mathcal{X}$，其中度量空间 $\mathcal{X}$ 配备了 Borel $\sigma$-代数。$X$ 在 $\mathcal{X}$ 上的分布 $P_X$，可由底层概率空间 $(\Omega,\mathcal{F},P)$ 中的测度 $P$ 导出。我们不必始终考虑这个底层空间，通常可以直接从随机变量的分布出发，假定随机试验直接给出从该分布抽取的数值。}
统计学和机器学习几乎都以独立同分布的数据为基础。然而，实际中的独立同分布假设可能以多种方式被破坏，例如发生\textbf{分布偏移（distribution shift）}，或系统受到\textbf{干预（intervention）}。我们将在后文看到，其中有些情形与因果性有着密切联系。

我们可能关心 $P_{X,Y}$ 的某些性质，例如：
\begin{enumerate}[label=(\roman*)]
 \item 给定输入时输出的期望，即 $f(x)=\mathbb{E}[Y\mid X=x]$。这称为\textbf{回归（regression）}，其中通常有 $\mathcal{Y}=\mathbb{R}$。
 \item 将每个 $x$ 归入概率较大类别的\textbf{二分类器（binary classifier）}，即
 \[
   f(x)=\argmax_{y\in\mathcal{Y}}P(Y=y\mid X=x),
   \qquad \mathcal{Y}=\{\pm1\}.
 \]
 \item $P_{X,Y}$ 的密度 $p_{X,Y}$，假定这一密度存在。
\end{enumerate}
在实际应用中，我们希望从有限数据集中估计这些性质。也就是说，我们依据样本 (1.1) 进行估计；等价地，也可以依据\textbf{经验分布（empirical distribution）} $P^n_{X,Y}$ 进行估计，这一分布在每个观测值上放置等权重的点质量。

这构成了一个\textbf{逆问题（inverse problem）}：我们希望估计某个无法直接观测的对象——这里是数据背后的分布——的某种性质，而所依据的观测值，是对这个对象施加某种操作后得到的。在这里，这一操作就是从未知分布中抽样。

\section{学习理论}
\label{sec:learning-theory}

现在，仿照从 $P_{X,Y}$ 得到 $f$ 的做法，假设我们用经验分布来构造经验估计 $f^n$。这实际上是一个\textbf{不适定问题（ill-posed problem）}［例如 Vapnik, 1998］，因为对样本 $(x_1,y_1),\ldots,(x_n,y_n)$ 中没有出现过的任何 $x$，相应的条件期望都没有定义。当然，我们可以先在观测样本上定义函数 $f$，再按照某种固定规则将其延拓到样本之外，例如在样本之外令 $f$ 等于 $+1$，或者选择一个连续的分段线性函数。然而，无论选择哪一种规则，输入——即经验分布——的微小变化，都可能引起输出的巨大变化。无论观测数量有多大，经验分布通常都不能完美地逼近真实分布；这种逼近中的微小误差，因而可能在估计中造成很大的误差。这意味着，如果不对经验估计 $f^n$ 所属的函数类作出额外假设，我们就无法保证这些估计能够在适当的意义下逼近最优目标 $f$。

统计学习理论借助\textbf{容量度量（capacity measure）}来形式化这些假设。如果所用函数类非常丰富，足以拟合几乎所有可以设想的数据集，那么能够拟合手头的数据并不令人意外。反之，如果事先将函数类的容量限制得很小，那么在所有可能的数据集中，只有少数数据集能够由该函数类中的某个函数来解释。在这种情况下，如果我们仍然能够解释手头的数据，就有理由相信，我们找到了蕴藏在数据中的某种规律。此时，对于未来从同一分布 $P_{X,Y}$ 抽取的数据，我们可以为解的准确性提供概率保证。

也可以从另一个角度理解这一点：我们所用的函数类融入了与观测数据背后规律相吻合的\textbf{先验知识（a priori knowledge）}，例如函数的光滑性。这类知识可以通过多种方式引入，而不同机器学习方法的差异之一，就在于如何处理这个问题。在贝叶斯方法中，我们为函数类和噪声模型指定先验分布；在正则化理论中，我们构造适当的正则化项，并将其纳入优化问题，以引导解的选择。

统计学习的困难，主要在于我们试图依据经验数据求解一个逆问题。如果已知完整的概率模型，这些问题便不复存在。讨论因果模型时，我们将看到，因果学习在某种意义上更加困难，因为它存在两个层面的不适定性。第一层是统计意义上的不适定性：无论有限样本有多大，它都不可能包含底层分布的全部信息。第二层不适定性则来自这样一个事实：即使完全知道观测分布，通常仍不能确定其背后的因果模型。

下面更详细地考察统计学习问题，并聚焦于\textbf{二类模式识别（binary pattern recognition）}，也就是二分类问题［例如 Vapnik, 1998］。此时，$\mathcal{Y}=\{\pm1\}$。我们希望根据从未知分布 $P_{X,Y}$ 独立同分布抽取的观测样本 (1.1)，学习函数 $f\colon\mathcal{X}\to\mathcal{Y}$。我们的目标是在某个函数类 $\mathcal{F}$ 中，使期望误差或\textbf{风险（risk）}\footnote{这里的“风险”概念来自统计学习理论［Vapnik, 1998］，其含义并不总与日常用法相同；它的源头可以追溯到统计决策理论［Wald, 1950；Ferguson, 1967；Berger, 1985］。在这一语境下，$f(x)$ 被理解为观测到 $x$ 后采取的行动，而损失函数衡量当真实状态为 $y$ 时，这一行动所造成的损失。}
\begin{equation}
 R[f]=\int \frac{1}{2}\lvert f(x)-y\rvert\,\mathrm{d}P_{X,Y}(x,y)
 \tag{1.2}
\end{equation}
最小。注意，这里是关于测度 $P_{X,Y}$ 的积分；如果 $P_{X,Y}$ 相对于 Lebesgue 测度存在密度 $p(x,y)$，那么这一积分就可以写成
\[
 \int \frac{1}{2}\lvert f(x)-y\rvert\,p(x,y)\,\mathrm{d}x\,\mathrm{d}y.
\]
由于 $P_{X,Y}$ 未知，我们无法计算式 (1.2)，更不用说将其最小化。因此，我们转而采用某种\textbf{归纳原则（induction principle）}，例如\textbf{经验风险最小化（empirical risk minimization）}。具体做法是在 $f\in\mathcal{F}$ 中，选择使训练误差或\textbf{经验风险（empirical risk）}
\begin{equation}
 R^n_{\mathrm{emp}}[f]
 =\frac{1}{n}\sum_{i=1}^{n}\frac{1}{2}\lvert f(x_i)-y_i\rvert
 \tag{1.3}
\end{equation}
最小的函数。

从渐近角度看，一个重要问题是，这样的学习过程是否具有\textbf{一致性（consistency）}。其基本含义是：随着 $n$ 趋于无穷，所产生的函数序列的风险，依概率收敛于给定函数类 $\mathcal{F}$ 中能够达到的最小风险。附录 A.3 将说明，只有当函数类足够“小”时，这一点才可能成立。\textbf{Vapnik--Chervonenkis（VC）维}［Vapnik, 1998］是度量函数类容量或大小的一种方式。它还能帮助我们得到有限样本保证：以较高概率，风险 (1.2) 不超过经验风险再加上一个随函数类 $\mathcal{F}$ 的大小而增大的项。

这一理论与已有的\textbf{普适一致性（universal consistency）}结果并不矛盾。普适一致性是指，学习算法的风险收敛于所有函数中能够达到的最低风险。确实存在具有普适一致性的学习算法，例如最近邻分类器和支持向量机［Devroye et al., 1996；Vapnik, 1998；Schölkopf and Smola, 2002；Steinwart and Christmann, 2008］。普适一致性在本质上告诉我们，在数据量趋于无穷的极限下，任何问题都可以学习；但它并不意味着，每个问题都能凭借有限数据得到很好的学习结果，因为还可能出现\textbf{慢收敛现象（slow rates）}。对于任意学习算法，总存在某些问题，使其学习收敛速率任意缓慢［Devroye et al., 1996］。不过，普适一致性确实表明，只要固定分布并收集足够多的数据，最终就能够任意逼近最低风险。

在实际应用中，机器学习系统近年来的成功似乎表明，有时我们确实已经进入上述渐近情形，并且往往取得了出色结果。人们投入了大量思考，力求设计出尽可能高效利用数据的方法，以便从给定数据集中得到尽可能好的结果；同时，也投入了大量精力构建可用于训练这些方法的大规模数据集。然而，在所有这些情形下，一个关键前提是：底层分布在训练与测试之间不能因干预或其他变化而发生改变。本书将论证，如果只是用概率分布描述底层规律，而不引入额外结构，我们就缺乏恰当的手段来描述究竟哪些部分可能发生变化。
